% Cover letter — Research Scientist, Doctolab (Doctolib), Paris
% Compile: xelatex coverletter-doctolib.tex
\documentclass[11pt,a4paper]{article}
\usepackage{fontspec}
\usepackage[english]{babel}
\usepackage[top=1.6cm, bottom=1.6cm, left=2.0cm, right=2.0cm]{geometry}
\usepackage{amssymb}
\usepackage{xcolor}
\usepackage{enumitem}
\usepackage{titlesec}
\usepackage[hidelinks]{hyperref}
\usepackage{fontawesome5}
\setmainfont{texgyrepagella}[Extension=.otf, UprightFont=*-regular, BoldFont=*-bold, ItalicFont=*-italic, BoldItalicFont=*-bolditalic]
\newfontfamily\headingfont{texgyreheros}[Extension=.otf, UprightFont=*-regular, BoldFont=*-bold, ItalicFont=*-italic, BoldItalicFont=*-bolditalic]
\definecolor{accent}{RGB}{28,58,94}
\definecolor{ink}{RGB}{26,26,26}
\definecolor{muted}{RGB}{95,105,120}
\definecolor{rulegray}{RGB}{200,206,214}
\color{ink}
\pagestyle{empty}
\setlength{\parindent}{0pt}
\setlength{\parskip}{0.6ex}
\linespread{1.03}
\titleformat{\section}{\headingfont\bfseries\fontsize{10}{12}\selectfont\color{accent}}{}{0em}{\MakeUppercase}[\vspace{-0.9ex}{\color{rulegray}\rule{\linewidth}{0.7pt}}]
\titlespacing*{\section}{0pt}{2.0ex plus .3ex}{1.0ex}
\titleformat{\paragraph}[runin]{\headingfont\bfseries\small\color{accent}}{}{0em}{}
\titlespacing*{\paragraph}{0pt}{1.0ex}{0.6em}
\setlist{itemsep=0.25ex, topsep=0.4ex, leftmargin=1.3em, label={\textcolor{accent}{\scriptsize$\blacktriangleright$}}}
\newcommand{\letterhead}{%
{\headingfont\bfseries\fontsize{16}{19}\selectfont EL MEHDI HARESS}\\[0.4ex]
{\footnotesize\color{muted}Inria Sophia Antipolis \textemdash{} CALISTO team\quad
\faEnvelope[regular]~\href{mailto:el-mehdi.haress@inria.fr}{el-mehdi.haress@inria.fr}\quad
\faPhone*~+33 7 58 23 41 50\quad
\faGlobe~\href{https://elmehdiharess.org}{elmehdiharess.org}}\\[-0.6ex]
{\color{rulegray}\rule{\linewidth}{0.7pt}}\par\vspace{1.2ex}}

\begin{document}
\letterhead

\begin{flushright}\small Sophia Antipolis, \today\end{flushright}

\textbf{To:} Doctolab, Doctolib Clinical AI Research Lab\\
\textbf{Subject:} Application for the Research Scientist position (Paris)

Dear Hiring Team,

I am applying for the Research Scientist position at Doctolab. I am an applied mathematician (PhD,
Université Paris-Saclay; engineer, CentraleSupélec), currently a research engineer at Inria Sophia
Antipolis after a two-year research fellowship at the University of Leeds. My research is about one
question that sits at the heart of your lab's agenda: \textbf{how a state evolves over time under
uncertainty, and how to learn that evolution reliably from imperfect observations.}

Your description of a world model of patient health, a health state that evolves over time and responds
to care, learned from observational data, is close to the objects I work with every day. Concretely, I
would bring three things.

\begin{itemize}
\item \textbf{Temporal and dynamic modelling.} I study stochastic dynamical systems, including systems
with long memory, and I have proved stability results that quantify how a trajectory reacts when the
model or its parameters are perturbed (\emph{Stochastic Processes and their Applications}, 2026). For a
world model, this is the question of when a prediction can still be trusted far from the training data.
\item \textbf{Calibration and inference from discrete, noisy data.} I have designed estimators that
recover several parameters jointly (drift, noise level, memory) from discretely sampled trajectories,
with convergence guarantees, and released the code. Irregular, partial observation of a continuous
process is exactly the situation of longitudinal patient records.
\item \textbf{Robustness and risk-aware learning.} At Osaka University I worked on learning algorithms
that optimise risk measures rather than averages, with guarantees under heavy-tailed losses (two
\emph{AISTATS} papers). In a clinical setting, rare outcomes are often the ones that matter most.
\end{itemize}

I have driven my own research from the choice of question to publication, and I am used to explaining
why a result matters to people outside my field: I taught probability, statistics and stochastic finance
at CentraleSupélec and in Leeds, and I now work with engineers and industrial partners at Inria.

I want to be transparent about where my profile is thinner. My experience with deep learning frameworks
and with large, messy real-world data sets is more limited than my background in modelling and theory,
and I have not yet worked with health data. These are the skills I most want to build, and I learn
fast in hands-on settings: I taught myself to design and ship full web applications in TypeScript and
React. What I can offer in return is a rigorous view of dynamics, calibration and uncertainty, which is
the part of a clinical world model that is hardest to get right.

The fact that Doctolab is in its founding phase, and that results reach practitioners and patients at
scale, is what makes this position particularly attractive to me. I would be glad to discuss how my work
could contribute to your research agenda.

Yours sincerely,\\[0.8ex]
El Mehdi Haress

\end{document}
